\documentclass[11pt]{article}
\usepackage[a4paper,margin=21mm]{geometry}
\usepackage{fontspec}
\setmainfont{Latin Modern Roman}
\newfontfamily\CJKfont{Noto Serif CJK SC}
\XeTeXlinebreaklocale "zh"
\XeTeXlinebreakskip=0pt plus 1pt
\usepackage{amsmath,amssymb,amsthm,amscd,graphicx,color}
\usepackage{xurl}
\usepackage[unicode,hidelinks]{hyperref}
\allowdisplaybreaks
\setlength\emergencystretch{3em}
\numberwithin{equation}{section}
\usepackage{lineno}
\usepackage{cancel}
	
\newcommand{\Z}{\mathbb Z}

\newcommand{\C}{\mathbb C}
\newcommand{\Q}{\mathbb Q}

\newcommand{\g}{\mathfrak g}

\newcommand{\F}{\mathcal F}
\newcommand{\W}{\mathcal W}

\newcommand{\hg}{\hat \g}

\newcommand{\bea}{\begin{eqnarray}}
\newcommand{\eea}{\end{eqnarray}}

\newtheorem{Theorem}{Theorem}[section]
\newtheorem*{Theorem*}{Theorem}
\newtheorem{Corollary}[Theorem]{Corollary}
\newtheorem{Lemma}[Theorem]{Lemma}
\newtheorem{Proposition}[Theorem]{Proposition}
\newtheorem{conj}[Theorem]{Conjecture}
\newtheorem{Claim}[Theorem]{Claim}

 { \theoremstyle{definition}
\newtheorem{Definition}[Theorem]{Definition}
\newtheorem{Note}[Theorem]{Note}
\newtheorem{Example}[Theorem]{Example}
\newtheorem{Remark}[Theorem]{Remark}
}

\newcommand{\DR}[1]{\textcolor{red}{#1}}
\newfontfamily\nativebbone{DejaVu Math TeX Gyre}
\ExplSyntaxOn
\NewDocumentCommand{\mathbbm}{m}{\str_if_eq:nnTF{#1}{1}{\mathord{\text{\nativebbone\char"1D7D9}}}{\PackageError{nativecompat}{Unsupported~mathbbm~argument}{Only~verified~digit~one~is~accepted}}}
\ExplSyntaxOff
\makeatletter\newlabel{KS_gen}{{3}{}{Original source locator}{}{}}\makeatother
\begin{document}
\begin{center}\Large Kazama–Suzuki duality at level minus one\end{center}
\noindent\textbf{Stable ID:} 561\quad\textbf{Author:} Dražen Adamović; Ana Kontrec\par
\noindent\textbf{Work:} On Kazama–Suzuki Duality between W\_k(sl\_4, f\_sub) and N=2 Superconformal Vertex Algebra\par
\noindent\textbf{Source:} \url{https://sigma-journal.com/2025/036/}\par
\noindent\textbf{Version:} Official SIGMA 21 (2025) PDF and native TeX; acquired 2026-09-30; exact bytes pinned by SHA256\par
\noindent\textbf{Rights:} CC BY 4.0. Authors retain copyright. \url{https://creativecommons.org/licenses/by/4.0/}. Reuse and typesetting adaptation; original language and mathematical content preserved.\par
\noindent\textbf{Difficulty (prior selection preserved):} {\CJKfont H2: Mutually inverse lattice realizations require proving simplicity and identifying the maximal ideal, not just verifying displayed operator products. The author’s spectral-flow/simple-current decomposition reconstructs every generator and proves the two commutants agree, giving a nonroutine quotient-identification core.}\par
\medskip\noindent\textbf{Editorial boundary note (not author text):} Complete original Theorem4.6, with both author embedding/OPE calculations, simplicity and simple-current decomposition arguments, required preceding singular-vector proof, lattice/spectral-flow/parafermion definitions, and the actual deferred OPE appendix. The final theorem records the preceding construction rather than giving a separate proof paragraph. Claim4.5 uses the stated analogy to Claim4.2, whose full proof is retained; no assistant proof is supplied. Later module-classification developments and the other-level classification are excluded. A narrowly guarded typography-only resolver preserves literal mathbbm calls for the verified vacuum/identity digit1 and renders the double-struck glyph using the named installed DejaVu Math TeX Gyre font; this is an explicit portable font dependency. All other arguments are rejected.\par
\noindent\textbf{External original reference KS\_gen:} Original Section3, source412–521, classifies the possible level pairs. Referenced in the section-opening recap only; the selected existence proof for k=-1 includes its own singular-vector and embedding constructions.\par
\medskip\hrule\medskip\noindent\textbf{Author text begins below.}\par
\setcounter{section}{1}\setcounter{equation}{0}
% Original source lines 103--110
Assume that $U$, $V$ are vertex (super)algebras. We say that $V$ is the \emph{Kazama--Suzuki dual} of~$U$ if there exist injective homomorphisms of vertex superalgebras
\begin{equation} \label{KSdef}
	\varphi_1 \colon\ V \rightarrow U \otimes \F_1, \qquad \varphi_2 \colon\ U \rightarrow V \otimes \F_{-1},
\end{equation}
so that
$V \cong \operatorname{Com} ( M_{H_1} (1) , U \otimes \F_1)$, $ U \cong \operatorname{Com} ( M_{H_2} (1) , V \otimes \F_{-1})$,
where $M_{H_1} (1) $ (resp.\ $ M_{H_2} (1)$) is a rank one Heisenberg vertex subalgebra of
$U\otimes \F_1$ (resp.\ $V\otimes \F_{-1}$) and $\F_{\pm 1}$ are lattice vertex superalgebras associated to rank one lattice $L= {\Z}\varphi^{\pm} $, $\big\langle \varphi^{\pm}, \varphi^{\pm} \big\rangle = \pm 1$ (cf.\ Section \ref{F_{-1}}).

\setcounter{section}{1}
% Original source lines 192--383
{\bf Setup.}
\begin{itemize}\itemsep=0pt
	\item We adopt the following notation for the vertex operator corresponding to the state $a$
\[
Y (a, z) = \sum _{i \in \mathbb{Z}} a_{(i)}z^{-i-1}.
\]
\item The Zhu algebra associated to the vertex operator algebra $V$ with the Virasoro vector $\omega$ will be denoted with $A_{\omega}(V)$.
\item Let $\F_{\pm 1}$ be the lattice vertex {super}algebras associated to the lattice ${\mathbb Z}\sqrt{\pm 1}$ defined in Section \ref{F_{-1}}.
\item Let $M_{\alpha}(1)$ denote the rank one Heisenberg vertex algebra generated by a Heisenberg vector~$\alpha$.
\item Let $\mathcal{W}^k(\mathfrak{sl}_n, f_{\rm sub})$ be the (universal) subregular $\mathcal{W}$-algebra corresponding to the subregular nilpotent element $f_{\rm sub}$ and $\mathfrak g = \mathfrak{sl}_n$, and $\mathcal{W}_k(\mathfrak{sl}_n, f_{\rm sub})$ its simple quotient.
\item $L_s(\mathfrak{sl}_2)$ is affine Lie algebra of $\mathfrak{sl}_2$ and $N_s(\mathfrak{sl}_2)$ its parafermion (coset) subalgebra.
\item The universal $N=2$ superconformal vertex algebra of central charge $c$ is denoted by $V^{N=2}_c$ and $L^{N=2}_{c}$ is its simple quotient.
\end{itemize}

\section{Preliminaries}	

\subsection[Lattice vertex superalgebras F\_1]{ Lattice vertex superalgebras $\boldsymbol{\mathcal F_{\pm 1}}$} \label{F_{-1}}

Consider a rank one lattice $L= {\Z}\varphi^{\pm} $, $\big\langle \varphi^{\pm}, \varphi^{\pm} \big\rangle = \pm 1$. Let $\F_{1}$ (resp.\ $\F_{-1}$) be the associated vertex algebra. These vertex superalgebras are used for the construction of the inverse of Kazama--Suzuki functor in the context of duality between affine $\widehat{sl}_2$ and $N=2$ superconformal algebra (cf.\ \cite{A-IMRN,A-2001,FST}).

As a vector space $\F_{\pm 1} = {\C}[L] \otimes M_{\varphi^{\pm}} (1)$, where ${\C}[L]$ is a group algebra of $L$, and $M_{\varphi^{\pm}} (1)$ the Heisenberg vertex algebra generated by the Heisenberg field $\varphi^{\pm}(z) = \sum_{n \in {\Z} } \varphi^{\pm}(n) z^{-n-1}$ such that
\[
 \bigl[\varphi^{\pm}(n), \varphi^{\pm}(m)\bigr] = \pm n \delta_{n+m, 0}.
 \]

The vertex algebras $\F_{\pm 1}$ are weakly generated by \smash{$e^{ \varphi^{\pm}}$} and \smash{$e^{-\varphi^{\pm}}$}.
Moreover, $\F_{\pm 1}$ is a simple vertex superalgebra and a completely reducible $M_{\varphi^{\pm}} (1)$-module isomorphic to
\[
\F_{\pm 1} \cong \bigoplus_{m \in \mathbb{Z}} \F_{\pm 1}^m ,
\]
where $ \F_{\pm 1}^m$ is an irreducible $M_{\varphi^{\pm}} (1)$-module generated by \smash{$e^{m\varphi^{\pm}}$}.


\subsection[N=2 superconformal algebra]{$\boldsymbol{ N=2}$ superconformal algebra}

The universal $N=2$ superconformal vertex algebra $V^{N=2}_{c}$ is generated by even fields $T(z)$, $H(z)$ and odd fields $E(z)$, $F(z)$
satisfying the following OPEs
\begin{gather*}
 T(z)T(w) \sim \frac{\frac{c}{2}}{(z-w)^{4}}+\frac{2T(w)}{(z-w)^{2}} + \frac{\partial T(w)}{z-w}, \qquad H(x)H(y) \sim \frac{ \frac{c}{3}}{(z-w)^{2}}, \qquad \\
 T(z)H(w) \sim \frac{H(w)}{ (z-w)^{2}}+ \frac{\partial H(w)}{z-w}, \\
 T(z)E(w) \sim \frac{\frac{3}{2} E(w)}{(z-w)^{2}} + \frac{\partial E(w)}{z-w}, \qquad T(z)F(w) \sim \frac{\frac{3}{2} F(w)}{(z-w)^{2}} + \frac{\partial F(w)}{z-w}, \\
	E(z)F(w) \sim \frac{\frac{2c}{3}}{(z-w)^{3}}+\frac{2H(w)}{(z-w)^{2}} + \frac{2T(w) + \partial H(w)}{z-w}, \\
	F(z)E(w) \sim \frac{\frac{2c}{3}}{(z-w)^{3}}-\frac{2H(w)}{(z-w)^{2}} + \frac{2T(w) - \partial H(w)}{z-w}, \\
	H(z)E(w) \sim \frac{E(w)}{z-w}, \qquad H(z)F(w) \sim -\frac{F(w)}{z-w}, \qquad E(z)E(w) \sim 0, \qquad F(z)F(w) \sim 0. 	
\end{gather*}

Set
\begin{gather*}
 T(z) = \sum _{i \in {\Z}} T(i) z^{-i-2}, \qquad H(z) = \sum _{i \in {\Z}} H(i) z^{-i-1},\\
 E(z) = \sum _{i \in {\Z}} E_{(i)} z^{-i-1} = \sum _{i \in {\Z}} E\left(i+\frac{1}{2}\right)z^{-i-2},\\
 F(z) = \sum _{i \in {\Z}} F_{(i)} z^{-i-1} = \sum _{i \in {\Z}} F\left(i+\frac{1}{2}\right)z^{-i-2}.
 \end{gather*}
 Then the components of these fields satisfy the commutation relation for the $N=2$ superconformal algebra with basis $\{T(n), H(n), E(r), F(r)\}$, $n \in \mathbb{Z}$, $r \in \frac{1}{2}+\mathbb{Z}$,
\begin{gather*}
	[T(m), T(n)] = (m-n)T(m+n) + \frac{c}{12}\bigl(m^3 - m\bigr)\delta_{m+n,0}, \qquad [H(m), H(n)] = \frac{c}{3}m\delta_{m+n,0}, \\
	[T(m), H(n)] = -nH(n+m), \qquad [ T(m) , E(r) ] = \left(\frac{1}{2} m-r \right)E(m+r), \\
 [ T(m) , F(r) ] = \left(\frac{1}{2} m-r \right)F(m+r), \\
	[ H(m) ,E(r)] = E(m+r), \qquad [ H(m) , F(r)] = -F(m+r), \\
	\{E(r), F(s) \} = 2 T({r + s}) + ( r - s )H(r + s) + \frac{c}{3}\left( r^2 - \frac{1}{4} \right)\delta_{r+s,0}, \\
	\{E(r), E(s) \} = \{F(r), F(s) \} = 0.
\end{gather*}

 Let $L^{N=2}_{c}$ be the simple quotient of $V^{N=2}_{c}$.

\subsubsection{Spectral flow automorphisms}
The $N=2$ superconformal algebra $V^{N=2}_{c}$ admits a family of spectral flow automorphisms $\sigma^\ell$, $\ell \in \mathbb{Z}$, given by
\begin{align*}
\begin{aligned}
	& \sigma^\ell(T(n))= T(n) -\ell H(n) + \frac{1}{6}\ell ^2\delta_{n,0}c\mathbbm{1}, \qquad \sigma^l(H(n)) = H(n) + \frac{1}{3}\ell \delta_{n,0}c\mathbbm{1},  \\
	& \sigma^\ell(E(r))= E({r-\ell}), \qquad \sigma^\ell(F(r)) = F({r+\ell}), \qquad \sigma^\ell(\mathbbm{1}) = \mathbbm{1}.
\end{aligned}
\end{align*}



Define
\[
 \Delta (-H,z) = z^{-H(0)}\exp \left ( \sum _{k=1} ^{\infty} (-1)^{k+1}\frac{-H(k) }{kz^k}\right ).
 \]

For any $V^{N=2}_{c}$-module $(M, Y_M(\cdot, z))$ and $n \in \mathbb{Z}$,
$\sigma^ n (M)$ is again a $L^{N=2}_{c}$-module with vertex operator structure given by
$ Y_{\sigma ^n (M) } ( \cdot , z)) := Y_{ M } (\Delta ( -n H, z) \cdot , z))$.

\subsubsection{Twisted highest weight conditions}
Let us define a new Virasoro vector $\widetilde T = T + \partial H$. Then $E(z)$ is primary field of conformal weight~$1/2$ and $F(z)$ of conformal weight $5/2$ with respect to $\widetilde T$.
 With respect to the new grading operator $\widetilde T(0)$ we shall also need to introduce twisted highest weight conditions.
 Let~${(h,q) \in {\C}^2}$. Let $M_c[h,q]$ be the Verma module of the $N=2$ superconformal algebra of central charge $c$ generated by the (twisted) highest weight vector $v_{h,q}$ such that for $n \in {\Z}_{\ge 0}$ (cf.\ \cite{FST}):
\begin{gather*}
 H(n) v_{h,q} =\delta_{n, 0} h v_{h,q}, \qquad T(n) v_{h,q} = \delta_{n,0} q v_{h,q}, \\ E\left(n-\frac{1}{2}\right) v_{h,q} = F\left(n+\frac{3}{2}\right) v_{h,q} =0.
 \end{gather*}
 Let $L_c[h,q]$ be the irreducible quotient of $M_c[h,q]$. Then $L_c[h,q]$ is an irreducible $V_c^{N=2}$-module.

\subsubsection{Parafermionic subalgebras}
\label{paraferm-1}
Let $ \mathcal N_c = \operatorname{Com} \bigl(M_H(1) , L^{N=2}_{c}\bigr) $ be the parafermionic subalgebra of $L^{N=2}_{c}$. Note that using the Kazama--Suzuki duality between $L^{N=2} _c$ and $L_k(\mathfrak{sl}_2)$ (cf.\ \cite{A-IMRN, FST}), we have that
$\mathcal N_c = \mathcal N_s (\mathfrak{sl}_2)$, where $c =\frac{3s}{s+2}$ and
$\mathcal N_s (\mathfrak{sl}_2)$ is the parafermion vertex subalgebra of $L_s(\mathfrak{sl}_2)$ (cf.\ \cite{DLY}). If $c \notin \{ 0, 1, \frac{3}{2}, -6, -9\}$, then the vertex algebra $ \mathcal N_c $ is (weakly) generated by the fields $ T^{\perp}$ and~$W^{N=2}_{c}$ (cf.\ \cite{BEHHH,DLY}), where
\begin{gather*}
	 T^{\perp} = T - \frac{3}{2c}{:}HH{:}, \qquad
	W^{N=2}_{c} = \nu \left({:} E F{:} - \partial T -\frac{6}{c}{:} T H{:} - \frac{c-9}{3c}\partial^2H + \frac{6}{c^2}{:}H^3{:} \right),
\end{gather*}
and $\nu \in \mathbb{C}$ is a normalization factor. In particular, $\mathcal N_{c=-15}$ is weakly generated by $T^{\perp}$ and $W^{N=2}_{c=-15}$.

 For other central charges, we have the following:
\begin{itemize}\itemsep=0pt
\item If $c\in \{ 0, 1\}$, then $\mathcal N_c = {\C} {\bf 1}$ (cf.\ \cite{DLY}).
\item $N_{c=3/2}$ is the simple Virasoro vertex algebra of central charge $c=\frac{1}{2}$.
 \item \smash{$\mathcal N_{c=-6}=\mathcal N_{-4/3}(\mathfrak{sl}_2)$} is the singlet vertex algebra $\mathcal M(3)$ generated by the Virasoro field and primary field of conformal weight $5$ (cf.\ \cite{A-2005}).
 \item \smash{$\mathcal N_{c=-9} =\mathcal N_{-3/2} (\mathfrak{sl}_2)$} is a direct sum of \smash{$\mathcal W_{-5/2} (\mathfrak{sl}_3, f_{\rm pr})$}-modules (cf.\ \cite{AMW}). As a vertex algebra it is generated by \smash{$\mathcal W_{-5/2} (\mathfrak{sl}_3, f_{\rm pr})$} and its simple module of conformal weight $4$.
\end{itemize}

 Let $A_{\omega}(V)$ be the Zhu algebra associated to the VOA $V$ with the Virasoro vector $\omega$, and let~$[v]$ be the image of $v \in V $ under the mapping $V \mapsto A_{\omega}(V)$.

 One can show that the Zhu's algebra $A\bigl(V^{N=2}_{c}\bigr)$ (cf.\ \cite{A-IMRN}) is isomorphic to ${\mathbb C}[x,y]$; i.e., there is isomorphism $f \colon A\bigl(V^{N=2}_{c}\bigr) \rightarrow {\mathbb C}[x_1,y_1]$ such that
$[H({-1}){\mathbbm 1}] \mapsto x_1, \ [T({-2}{)\mathbbm 1}] \mapsto y_1$.

\subsection{Kazama--Suzuki duality}
In this subsection, we will define a duality of vertex algebras which is motivated by the duality between $N=2$ superconformal vertex algebra and affine vertex algebra $L_k(\mathfrak{sl}_2)$.

Recall first that if $S$ is a vertex subalgebra of $V$, we have the commutant subalgebra of $V$ (cf.\ \cite{LL})
\[
\operatorname{Com} (S, V):= \{ v \in V \mid a_n v = 0,\, \forall a \in S,\, \forall n \in {\Z}_{\ge 0}\}.
\]

Assume that $U$, $V$ are vertex superalgebras. We say that $V$ is the Kazama--Suzuki dual of~$U$ if there exist injective homomorphisms of vertex superalgebras
$
 \varphi_1 \colon V \rightarrow U \otimes \F_1$, $ \varphi_2 \colon U \rightarrow V \otimes \F_{-1}$,
so that
$
V \cong \operatorname{Com} \left( M_{H_1} (1), U \otimes \F_1\right)$, $ U \cong \operatorname{Com} \left( M_{H_2} (1), V \otimes \F_{-1}\right)$,
where $M_{H_1} (1) $ (resp.\ $ M_{H_2} (1)$) is a rank one Heisenberg vertex subalgebra of
$U\otimes \F_1$ (resp.\ $V\otimes \F_{-1}$), and $\F_{\pm 1}$ are lattice vertex superalgebras defined in Section~\ref{F_{-1}}.

\subsection[Subregular W-algebras W\^{}k(sl\_4, f\_sub]{Subregular $\boldsymbol{\mathcal{W}}$-algebras $\boldsymbol{\mathcal{W}^k(\mathfrak{sl}_4, f_{\rm sub})}$}	\label{subreg_ope}

Let $\mathfrak{g}$ be a simple finite-dimensional Lie algebra and $k \in \mathbb{C}$. Let $\mathcal{W}^k(\mathfrak{g}, f_{\rm sub})$ be the (universal) subregular $\mathcal{W}$-algebra corresponding to the subregular nilpotent element $f_{\rm sub}$ (\cite{KW}).

In the case where $\mathfrak g = \mathfrak{sl}_n$, it was proved in \cite{Gen} that $\mathcal{W}^k(\mathfrak{g}, f_{\rm sub})$ is isomorphic to the Feigin--Semikhatov algebra \smash{$W_n^{(2)}$} (cf.\ \cite{FS}), using a certain free field realization of $\mathcal{W}^k(\mathfrak{g}, f_{\rm sub})$.

In this paper, we will only consider the case $\mathfrak g = \mathfrak{sl}_4$. The vertex algebra $\mathcal{W}^{k}(\mathfrak{sl}_4, f_{\rm sub})$ has central charge
\[
c_k = -\frac{(3k+8)(8k+17)}{(k+4) },
\]
 and it is freely generated by fields $J(z)$, $L(z)$, $G^+ (z)$, $G^- (z)$, $W(z)$. The explicit OPEs are known (see \cite{CL18,FS}) and are written in Appendix~\ref{appendix-b}. In this paper, we shall assume that the gradation on~$\mathcal{W}^{k}(\mathfrak{sl}_4, f_{\rm sub})$ is defined by the shifted Virasoro field $\widetilde L = L + \partial J$. Then $G^+$ (resp.~$G^-$) is a~primary field of conformal weight $1$ (resp.~$3$). We set
\begin{gather*}
J(z) = \sum _{n \in {\Z}} J(n) z^{-n-1}, \qquad L(z) = \sum _{n \in {\Z}} L(n) z^{-n-2}, \qquad G^+ (z) = \sum _{n \in {\Z}} G^+(n) z^{-n-1}, \\ G^- (z) = \sum _{n \in {\Z}} G^-(n) z^{-n-3}, \qquad W (z) = \sum _{n \in {\Z}} W(n) z^{-n-3}.
\end{gather*}	
The Zhu algebra $A\bigl(\mathcal{W}^{k}(\mathfrak{sl}_4, f_{\rm sub})\bigr)$ is generated by
\[
 E=[G^+],\qquad F=[G^- ], \qquad X=[J ], \qquad Y =[L], \qquad Z=[W].
 \]

The subregular $\mathcal{W}$-algebra $\mathcal{W}^k(\mathfrak{sl}_4, f_{\rm sub})$ admits a family of spectral flow automorphisms $\psi^m$, $m \in \mathbb{Z}$, given by
\begin{gather*}
 \psi^m (J(n)) = J(n)- { \frac{m (3 k +8)}{4} }
	 \delta _{n,0}\mathbbm{1} , \qquad \psi^m ( L(n)) = L(n) - mJ(n) + \frac{m^2 { (3k+8)}}{{8} }\delta _{n,0}\mathbbm{1}, \\
 \psi^m (G^+(n)) = G^+(n-m), \qquad \psi^m (G^-(n)) = G^-(n+m).
 \end{gather*}

Define
 \[
 \Delta (-J,z) = z^{-J(0)}\exp \left ( \sum _{k=1} ^{\infty} (-1)^{k+1}\frac{-J(k)}{kz^k}\right ).
 \]

For any $\mathcal{W}^k(\mathfrak{sl}_4, f_{\rm sub})$-module $(M, Y_M(\cdot, z))$ and $n \in \mathbb{Z}$,
$\psi^m (M)$ is again a $\mathcal{W}^k(\mathfrak{sl}_4, f_{\rm sub})$-module with vertex operator structure given by
$ Y_{\psi^m (M) } ( \cdot , z)) := Y_{ M } (\Delta ( -m J, z) \cdot , z))$.

\subsection[Singular vectors in W\^{}k(sl\_n, f\_sub)]{Singular vectors in $\boldsymbol{\mathcal{W}^k(\mathfrak{sl}_n, f_{\rm sub})}$}
	
	Formulas for singular vectors in $\mathcal{W}^k(\mathfrak{sl}_n, f_{\rm sub})$ are not known in general; however, the following criterion from \cite{Feh} tells us when $(G^+)^s$ is singular.
	
	\begin{Proposition}[\cite{Feh}]\label{kriterij}
		The vector $(G^+)^s$, $s>0$, is singular in $\mathcal{W}^k(\mathfrak{sl}_n, f_{\rm sub})$ if and only if~${
		 	i(k+n-1)=s}
$
 for some $i \in \{1,\dots ,n-1 \}$.
	\end{Proposition}

Using the above criterion, it is easy to observe the following.
	\begin{Lemma} \label{sing_G^+}\qquad
	\begin{itemize}\itemsep=0pt
		\item[$(i)$] For $k=-n+3$, $(G^+)^2$ is singular in $\mathcal{W}^k(\mathfrak{sl}_n, f_{\rm sub})$ but $G^+$ is not.
		\item[$(ii)$] For $k=-n+4$, $(G^+)^3$ is singular in $\mathcal{W}^k(\mathfrak{sl}_n, f_{\rm sub})$ but $(G^+)^2$ is not.	
	\end{itemize}
	\end{Lemma}
	
	\begin{proof}
		(i) Assume that $(G^+)^2$ is singular in $\mathcal{W}^k(\mathfrak{sl}_n), f_{\rm sub})$ but $G^+$ is not. Then from Proposition \ref{kriterij} it follows that there exists $i \in \{1,\dots ,n-1 \}$ such that $ i(k+n-1)=2 $ and there does not exist $i \in \{1,\dots ,n-1 \}$ such that $ i(k+n-1)=1$. Let now $k=-n+3$. Then $i=1$ is solution of the first equation, and the second equation does not have solutions. This proves (i).
		Analogous reasoning yields (ii).
	\end{proof}

\setcounter{section}{3}
% Original source lines 523--798
\section[The duality of W\_-1(sl\_4, f\_sub) and L\^{}N=2\_c=-15]{The duality of $\boldsymbol{\mathcal{W}_{-1}(\mathfrak{sl}_4, f_{\rm sub})}$ and $\boldsymbol{L^{N=2}_{c=-15}}$} \label{duality}


As we have seen in Section \ref{KS_gen}, Kazama--Suzuki duality between $L^{N=2}_c$ and $\mathcal W_k(\mathfrak{sl}_4, f_{\rm sub})$ can only occur if $k=-1$ and $c=-15$ or $k=-\frac{7}{3}$ and $c=1$.
In this section, we consider the case~${k=-1}$. Then the central charge of $\mathcal{W}^{k}(\mathfrak{sl}_4, f_{\rm sub})$ is $c = -15$. According to Lemma \ref{sing_G^+}, \smash{$(G^+)^2$} is singular vector in $\mathcal{W}^{k}(\mathfrak{sl}_4, f_{\rm sub})$, hence \smash{$(G^+)^2 = 0$} in $\mathcal{W}_{k}(\mathfrak{sl}_4, f_{\rm sub})$.

 We will show that $\mathcal{W}_{-1}(\mathfrak{sl}_4, f_{\rm sub})$ is the Kazama--Suzuki dual of the $N=2$ superconformal vertex algebra of central charge $-15$.



\subsection[Embedding of L\^{}N=2\_c=-15 into W\_-1(sl\_4, f\_ ub)]{Embedding of $\boldsymbol{ L^{N=2}_{c=-15}}$ into $\boldsymbol{\mathcal{W}_{-1}(\mathfrak{sl}_4, f_{\rm sub}) \otimes \F_{-1}}$}

First, we will show that the $N=2$ superconformal algebra $L^{N=2}_{c}$ for $c=-15$ can be realized as a subalgebra of $\mathcal{W}_{-1}(\mathfrak{sl}_4, f_{\rm sub}) \otimes \F_{-1}$, where $\F_{\pm 1}$ is the vertex algebra associated to the rank one lattice $L= {\mathbb{Z}}\varphi^{\pm}$, with $\big\langle \varphi^{\pm}, \varphi^{\pm} \big\rangle = \pm 1$ (cf.\ Section \ref{F_{-1}}).

 In this section, we shall use that fact that the universal superconformal $N=2$ vertex algebra~$V^{N=2}_{c}$ is simple for $c=-15$, since it is the Kazama--Suzuki dual of the universal affine vertex algebra $V^{k=-5/3}(\mathfrak{sl}_2)$ (which is also simple by \cite{GK}).

 Note that the maximal ideal in $\mathcal W^k(\mathfrak{sl}_4, f_{\rm sub})$ is invariant under the automorphism which maps $G^+$ to $G^-$. Since \smash{$(G^+)^2$} is a singular vector in $\mathcal{W}^{-1}(\mathfrak{sl}_4, f_{\rm sub}) $, we conclude that $(G^-)^2$ also belongs to the maximal ideal of $\mathcal{W}^{-1}(\mathfrak{sl}_4, f_{\rm sub}) $. Now we want to show that the maximal ideal is generated by these two vectors.
Let \smash{$I = \bigl\langle (G^+)^2, (G^-)^2 \bigr\rangle$} and define
\[
 \widetilde W = \mathcal{W}^{-1}(\mathfrak{sl}_4, f_{\rm sub})/I.
 \]

Let $H^{\perp} = J + \varphi^-$. Then for $n \ge 0$, we have
\[
H^{\perp}(n) H^{\perp} = \frac{1}{4} \delta_{n,1}\mathbbm{1}.
\]
Let $M_{H^{\perp}} (1)$ be the Heisenberg vertex algebra generated by $H^{\perp}$, and $M_{H^{\perp}} (1, s)$ the irreducible highest weight $M_{H^{\perp}} (1)$-module on which
$H^{\perp}(0) \equiv s \operatorname{Id}$.

\begin{Proposition} \label{ulaganje-1}\qquad
\begin{itemize}\itemsep=0pt
\item[$(1)$]There is a vertex algebra homomorphism $\Phi\colon L^{N=2}_{c=-15} \rightarrow \widetilde W \otimes \F_{-1}$ given by
	\begin{gather*}
		E = \frac{2}{3} G^+ \otimes e^{\varphi^-}, \qquad
		F = - G^- \otimes e^{-\varphi^-},\qquad
		H = - 4 J - 5 \varphi^-, \\
		T = L- 2{:}J J {:}-4{:}J\varphi^-{:} -\frac{5}{2}{:}(\varphi^-)^2{:}.
	\end{gather*}
	\item[$(2)$]
	Let $H^{\perp} = J + \varphi^-$. Then
\[
 \widetilde W \otimes \F_{-1} \cong \bigoplus_{n \in {\mathbb{Z}}} \sigma^n \bigl(L^{N=2}_{c=-15}\bigr) \otimes M_{H^\perp} (1, n).
 \]
	\item[$(3)$] We have
\[
\operatorname{Im} (\Phi) \cong \operatorname{Com} (M_{H^{\perp}} (1), \mathcal{W}_{-1}(\mathfrak{sl}_4, f_{\rm sub}) \otimes \F_{-1}) \cong L^{N=2}_{c=-15} .
\]
\item[$(4)$] $\widetilde W = \W_{-1} (\mathfrak{sl}_4, f_{\rm sub})$, i.e., $I$ is the maximal ideal in $\W^{-1}(\mathfrak{sl}_4, f_{\rm sub})$.
\end{itemize}
	
\end{Proposition}

\begin{proof}


It is easy to check that
\begin{gather*}
 H(0) E = E, \qquad H(0) F =-F, \qquad H(n) E = H(n) F =0, \qquad n >0, \\
 H(1) H = - 5 {\mathbbm 1} = \frac{c}{3} {\mathbbm 1},\qquad H(n) H =0, \qquad n >1.
\end{gather*}
 Let $L^{\perp} = L- \frac{2}{5} {:}J J{:}$ (cf.\ Appendix \ref{appendix-b}). We get that
\[
 T = L^{\perp} -\frac{1}{10} {:}H H{:}.
 \]
This implies that $T$ is a Virasoro vector of central charge $c=-15$. Clearly, we have that $H$ is a primary vector for $T$ of conformal weight $1$.


 Next we notice that $\bigl(G^\pm \bigr)^2 =0$ in $\widetilde W$. Direct calculation shows that
 \begin{gather*}
 E_{(0)} E = \frac{4}{9} {\bigl( G^+ \bigr)}^2 \otimes e^{2 \varphi^-} =0, \qquad F_{(0)}F = {( G^-)}^2 \otimes e^{-2 \varphi^-} =0, \\
 E_{(n)} E = F_{(n)} F =0, \qquad n\in {\mathbb Z}_{\ge 0},
\end{gather*}
and
\begin{gather*}
	E_{(2)} F = -\frac{2 (2 + k)(5 + 2k)(8 + 3k)}{3} {\mathbbm 1} =(-10) {\mathbbm 1} = \frac{2 c}{3} {\mathbbm 1}, \\
E_{(1)} F = -\frac{ 8(2 + k)(5 + 2k)}{3} J - \frac{2 (2 + k)(5 + 2k)(8 + 3k)}{3} \varphi^- = -8 J - 10 \varphi^- = 2 H, \\
 E_{(0)} F = - \frac{2 (2 + k)(5 + 2k)(8 + 3k)}{3}\cdot\frac{1}{2}\bigl({:}( \varphi^-)^2{:}+\partial\varphi^-\bigr)-\frac{ 8(2 + k)(5 + 2k)}{3} {:}J \varphi^-{:} \\
 \phantom{E_{(0)} F =}{}+\frac{2(k+2) \bigl((k+4)L-6{:}J^2{:}-2(2k+5) \partial J\bigr)}{3} = 2T + \partial H.
\end{gather*}

The above relations show that the even fields $H$, $ T$ and the odd fields $E$, $F$ satisfy the $\lambda$-bracket for the $N=2$ superconformal vertex algebra of central charge $c=-15$. Since \smash{$V^{N=2} _{c=-15}$} is simple, we conclude that $H$, $T$, $E$, $F$ generate a vertex subalgebra of $\widetilde W \otimes \F_{-1}$ isomorphic to~$L^{N=2}_{c=-15}$. This proves the assertion~(1).

 Let \smash{$ \overline W = \operatorname{Ker} _{ \widetilde W\otimes F_{-1}} H^{\perp}(0)$}. Then $\overline{W}$ is a vertex algebra which contains $ \operatorname{Im} (\Phi)\otimes M_{H^{\perp}} (1)$.



Let us prove the following claim.
\begin{Claim}\label{Claim1}
 $ \overline W $ is a simple vertex algebra generated by $\big\{ E,F,T,H, H^{\perp} \big\}$.
\end{Claim}

Let $U$ be the vertex subalgebra of $\overline W $ generated by $ \big\{E,F,T,H, H^{\perp} \big\}$. Then $U \cong \operatorname{Im} (\Phi) \otimes M_{H^{\perp}} (1)$. Since \smash{$\operatorname{Im} (\Phi) \cong L^{N=2}_{c=-15}$}, we conclude that
\smash{$U\cong L^{N=2}_{c=-15} \otimes M_{H^{\perp}} (1)$}, and it is therefore simple.

 For each $n \in {\Z}$, we consider the $U$-module $U^{(n)} = U. e^{n \varphi^-}$. One can show that
 $ U^{(n)} $ is isomorphic to the simple $U$-module obtained by the simple current construction
\[
 \bigl(U^{(n)}, Y^{(n)} (\cdot, z)\bigr): = (U , Y(\Delta(n \varphi^-,) \cdot, z)).
 \]


 Note that $ \varphi^- = -H - 4 H^{\perp}$, which implies that
 \[% \label{delta-alpha}
 \Delta(\varphi^-, z) = \Delta( -H, z) \Delta\bigl(- 4 H^{\perp}, z\bigr).
 \]
 \begin{itemize}\itemsep=0pt
\item For a $ \operatorname{Im} (\Phi)$-module $M$, by applying the operator $\Delta(-nH, z)$, we get the module $\sigma^{n}( M)$.
\item Applying the operator $ \Delta\bigl( -4n H^{\perp}, z\bigr)$ on
 $ M_{H^{\perp}}(1)$, we get the $ M_{H^{\perp}}(1)$-module
$ M_{H^{\perp}} (1, n )$.
\end{itemize}

We get
$
U^{(n)}= \sigma^{n}( \operatorname{Im} (\Phi)) \otimes M_{H^{\perp}} (1,n )$.
Note that $ H^{\perp} (0) \equiv n \operatorname{Id} $ on $U^{(n)}$.
 Using the construction of H. Li from \cite{Li-ext}, we get that
\[%\label{dec-u1}
 \mathcal U = \bigoplus_{n \in \mathbb{Z}} U^{(n)}
 \]
is a vertex algebra and hence a vertex subalgebra of $ \widetilde W \otimes \F_{-1}$. But it is not hard to see that $ \mathcal U$ contains all generators of $ \widetilde W \otimes \F_{-1}$.

 Indeed, $e^{\varphi^-}\in U^{(1)}$, $e^{-\varphi^-}\in U^{(-1)}$ and
 \[
 G^+ = \frac{3}{2} E_{(-2)}e^{-\varphi^-}, \qquad G^- =- F_{(-2)}e^{\varphi^-}.
 \]


Since $G^+$ and $G^-$ (weakly) generate $ \widetilde W$, it follows that
$
 \mathcal U = \widetilde W \otimes \F_{-1}$.






This proves that $ U= \overline W \cong \operatorname{Im} (\Phi) \otimes M_{H^{\perp}} (1) $. Therefore, $\overline W$ is a simple vertex algebra. This proves Claim~\ref{Claim1}.

Now we get
\[
 \widetilde W \otimes \F_{-1} = \bigoplus_{n \in \mathbb{Z}} U^{(n)} = \bigoplus_{n \in \mathbb{Z}} \sigma^n \bigl(L^{N=2}_{c=-15}\bigr) \otimes M_{H^\perp} (1, n).
 \]
This implies the assertions (2) and (3).

As the right-hand side of decomposition in (2) is simple, it follows that $\widetilde W $ must be simple and therefore isomorphic to the unique simple quotient $\mathcal{W}_{-1}(\mathfrak{sl}_4, f_{\rm sub})$. Hence $I$ is the maximal ideal in $\mathcal{W}^{-1}(\mathfrak{sl}_4, f_{\rm sub})$.
This proves assertion (4).
\end{proof}

 Let $M_H(1)$ be the Heisenberg subalgebra of $L^{N=2}_{c=-15}$ generated by the Heisenberg field $H(z)$ and $M_J(1)$ the Heisenberg subalgebra of $\mathcal{W}_{k}(\mathfrak{sl}_4, f_{\rm sub}) $ generated by the Heisenberg field $J(z)$.
From the proof of Proposition~\ref{ulaganje-1}, we have the following corollary.

	\begin{Corollary} \label{pf} $\operatorname{Com} \bigl(M_H(1) , L^{N=2}_{c=-15} \bigr) \cong \operatorname{Com} (M_J(1), \mathcal{W}_{-1}(\mathfrak{sl}_4, f_{\rm sub}) )$. \end{Corollary}
	
\subsection[Embedding of W\_k(sl\_4, f\_sub) into L\^{}N=2\_c=-15]{Embedding of $\boldsymbol{\mathcal{W}_{k}(\mathfrak{sl}_4, f_{\rm sub})}$ into $\boldsymbol{L^{N=2}_{c=-15} \otimes \F_{1}}$ }

It remains to show that $\mathcal{W}_{-1}(\mathfrak{sl}_4, f_{\rm sub}) $ can be realized as a subalgebra of the tensor product of the $N=2$ superconformal algebra $L^{N=2}_{c}$ for $c=-15$ and the lattice vertex algebra $\F_1$.

Let $\overline{H} = H - \varphi^+$. Then for $n \ge 0$ we have $\overline H(n) \overline H = -4 \delta_{n,1}\mathbbm{1}$.
Let $M_{\overline{H}} (1)$ be the Heisenberg vertex algebra generated by $\overline H$, and $M_{\overline{H}} (1, s)$ the irreducible highest weight $M_{\overline{H}} (1)$-module on which
$\overline H(0) \equiv s \operatorname{Id}$.





\begin{Proposition} \label{ulaganje-2}
\quad
\begin{itemize}\itemsep=0pt
	\item[$(1)$] There is a vertex algebra homomorphism $\Phi^{\rm inv} \colon \mathcal{W}^{-1}(\mathfrak{sl}_4, f_{\rm sub}) \rightarrow L^{N=2}_{c=-15} \otimes \F_{1}$ given by
	\begin{gather*}
		G^+ = E\otimes e^{\varphi^+}, \qquad
		G^- = \frac{3}{2} F \otimes e^{-\varphi^+}, \qquad
		J = - \frac{1}{4} H + \frac{5}{4} \varphi^+, \\
		L = T + \frac{1} {10} {:}H H{:}+ \frac{2}{5} {:}\left( - \frac{1}{4} H + \frac{5}{4} \varphi^+\right)^2{:}, \qquad
		W = -\frac{3}{2} W^{N=2}_{c=-15}.
	\end{gather*}
	
\item[$(2)$] $\operatorname{Im} \bigl(\Phi^{\rm inv}\bigr) $ is isomorphic to the simple vertex algebra $ \mathcal{W}_{-1}(\mathfrak{sl}_4, f_{\rm sub})$.

\item[$(3)$]
 As a $ \mathcal{W}_{-1}(\mathfrak{sl}_4, f_{\rm sub})\otimes M_{\overline{H}} (1)$-module
\[
L^{N=2}_{c=-15} \otimes \F_{1} \cong \bigoplus_{n\in {\Z} }
 \psi^{-n} ( \mathcal{W}_{-1}(\mathfrak{sl}_4, f_{\rm sub})) \otimes M_{\overline{H}} (1, n).
 \]
 \item[$(4)$] We have
\[
 \operatorname{Com} \bigl(M_{\overline{H}}(1), L^{N=2}_{c=-15} \otimes \F_{1}\bigr) \cong \mathcal{W}_{-1}(\mathfrak{sl}_4, f_{\rm sub}).
 \]
 \end{itemize}
	
\end{Proposition}
\begin{proof}
(1)
 First we notice that
$L = T^{\perp} + \frac{2}{5}{:}JJ{:}$,
which implies that $L$ is a Virasoro vector of central charge $c=-15$. Note that
$
 L^{\perp} = L- \frac{2}{5}{:}JJ{:} = T^{\perp}$.
Clearly, $J$ is a primary vector of conformal weight $1$ with respect to $L$.
Moreover, since
$L = T + \frac{2}{5}{:}JJ{:} + \frac{1}{10} {:}H H{:} $, we conclude that $G^{\pm}$ are primary vectors of conformal weight $2$.
Direct calculation shows that
\begin{gather*}
 G^+(3) G^- = \frac{3}{2}(-E_{(2)}F) = \frac{3}{2}\left(- \frac{2 c}{3}\right) \mathbbm{1} =15\mathbbm{1} =(2 + k)(5 + 2k)(8 + 3k)\mathbbm{1}, \\
 G^+(2) G^- =\frac{3}{2}\left(-E_{(1)}F- \varphi^+(-1) E_{(2)}F\right) \\
 \phantom{G^+(2) G^- }{} = \frac{3}{2}\left(-2H -\frac{2c}{3} \varphi^+\right) = 12\left( - \frac{1}{4} H + \frac{5}{4} \varphi^+ \right) =12J = 4(2 + k)(5 + 2k)J.
 \end{gather*}


 We should show that
\[ G^+(1) G^-
 = - 3 L + 6 {:} JJ {:} + 6 \partial J = -3 L^{\perp} + \frac{24}{5} J^2 + 6 \partial J.
 \]
 Indeed
\begin{align*}
G^+(1) G^- &= \frac{3}{2}\left(-E_{(0)}F-\varphi^+(-1) E_{(1)}F-\frac{1}{2}\bigl({:}\bigl(\varphi^+\bigr)^2{:}+\partial\varphi^+\bigr)E_{(2)}F \right) \\
&= \frac{3}{2}\left( -2 T -\partial H -2{:}H\varphi^+{:} - \frac{1}{2}\bigl({:}\bigl(\varphi^+\bigr)^2{:}+\partial\varphi^+ \bigr)\frac{2c}{3}\right)\mathbbm{1} \\
& = -3 L^{\perp}+ \frac{3}{10} {:}H^2{:}- \frac{3}{2} \partial H -3 {:}H \varphi^+ {:} + \frac{15}{2} \bigl({:}\bigl(\varphi^+\bigr)^2{:} +\partial \varphi^+\bigr) \\
& = -3 L^{\perp} + \frac{24}{5} {:}J ^2{:} + 6\partial J.
 \end{align*}

 Next, we need to show that
 \begin{align}
G^{+}(0) G^{-} ={}&(k+2)\left(W + \frac{8(11k+32)}{3(3k+8)^2}{:}J^3{:}-\frac{4(k+4)}{3k+8}{:}L J{:} + 6{:}\partial JJ{:}\right. - \frac{k+4}{2}\partial L \nonumber\\
 &+ \left. \frac{4(3k^2+17k+26)}{3(3k+8)}\partial^2J\right) \nonumber \\
={}& W+ \frac{32}{25}{:}J^3{:}-\frac{12}{5}{:}JL^{\perp}{:}+\frac{24}{5} {:}\partial JJ{:}-\frac{3}{2}\partial L^{\perp}+2\partial^2J.\label{GpGm-OPE}
\end{align}

We have that
 \begin{align}
G^+(0) G^-={}& \frac{3}{2}\left(-E_{(-1)}F -\varphi^+(-1) E_{(0)}F - E_{(1)}F\cdot \frac{1}{2}\bigl({:}\bigl(\varphi^+\bigr)^2{:}+\partial \varphi^+\bigr) \right. \nonumber\\
&\left. - E_{(2)}F\cdot \frac{1}{6}\bigl({:} \bigl(\varphi^+ \bigr)^3 {:}+3{:}\varphi ^+ \partial \varphi^+{:} + \partial ^2 \varphi^+\bigr) \right) \nonumber\\
={}& \frac{3}{2}\left(-E_{(-1)}F -\varphi^+(-1) \left(2L^{\perp} - \frac{1}{5}{:}H^2{:} +\partial H\right) - 2H\cdot\frac{1}{2}\bigl({:}\bigl(\varphi^+\bigr)^2{:}+\partial \varphi^+\bigr) \right. \nonumber\\
&\left.- \frac{2c}{3}\cdot \frac{1}{6}\bigl({:}\bigl(\varphi^+\bigr)^3{:} + 3{:}\varphi^+ \partial \varphi^+{:} + \partial ^2 \varphi^+\bigr) \right) .  \label{GpGm-rel}
\end{align}




As the parafermionic subalgebras of $\mathcal W_{-1} (\mathfrak{sl}_4, f_{\rm sub})$ and $L^{N=2}_{c=-15} $ coincide (cf.\ Corollary \ref{pf}), it follows that the field $W$ coincides (up to normalization) with the parafermionic generator~$W^{N=2}_{c=-15}$, that is,
\begin{align}
	W	&= \nu \left({:}E F{:} - \partial L^{\perp} - \frac{3}{2c}{:}H\partial H{:} -\frac{6}{c}{:}HL^{\perp}{:}- \frac{1}{3}\partial^2H + \frac{3}{c^2}{:}H^3{:} \right) \nonumber \\
	&= \nu \left({:}E F{:} - \partial L^{\perp}+ \frac{1}{5}{:}\partial H H {:} + \frac{2}{5}{:}HL^{\perp}{:} -\frac{1}{3}\partial^2H - \frac{1}{75}{:}H^3{:} \right). \label{W}
\end{align}
Setting $\nu = -\frac{3}{2}$ and substituting \eqref{W} into \eqref{GpGm-OPE}, we obtain that \eqref{GpGm-rel} = \eqref{GpGm-OPE}.
This proves the assertion (1).

(2) Let us prove that $\operatorname{Im} \bigl(\Phi^{\rm inv}\bigr) $ is simple. Let \smash{$\overline{W} = \operatorname{Ker}_{ L^{N=2}_{c=-15} \otimes \F_{1}} \overline{H} (0)$}. It is clear that $\overline{W}$ is a simple vertex algebra which contains $ \operatorname{Im} \bigl(\Phi^{\rm inv}\bigr) \otimes M_{\overline{H}} (1)$. The simplicity of $ \operatorname{Im} \bigl(\Phi^{\rm inv}\bigr) $ follows from the following claim.

\begin{Claim}\label{Claim2}
$\overline W$ is generated by $\big\{ G^+, G^-, J, L, W, \overline{H} \big\}$.
\end{Claim}

Proof of Claim~\ref{Claim2} is analogous to the proof of Claim~\ref{Claim1} in Proposition \ref{ulaganje-1}. Let $U$ be the vertex subalgebra of $\overline W$ generated by $\big\{ G^+, G^-, J, L^U, W, \overline{H} \big\}$. Then clearly $U \cong \operatorname{Im} \bigl(\Phi^{\rm inv}\bigr) \otimes M_{\overline{H}} (1)$.

Let $ U^{(n)} $ be the $U$-module obtained by the simple current construction
\[ \bigl(U^{(n)}, Y^{(n)} (\cdot, z)\bigr): = \bigl(U , Y\bigl(\Delta\bigl(n \varphi^+\bigr) \cdot, z\bigr)\bigr).
\]

As in Proposition \ref{ulaganje-1}, from the formula
\[% \label{delta-alpha}
\Delta\bigl(n\varphi^+, z\bigr) = \Delta( nJ, z) \Delta\left(\frac{n}{4} \overline{H}, z\right)
\]
we get
$U^{(n)}= \psi^{-n}\bigl( \operatorname{Im} \bigl(\Phi^{\rm inv}\bigr)\bigr) \otimes M_{\overline{H}} (1,n )$.


The rest of the proof follows analogously.
\end{proof}

We have proved the following.

\begin{Theorem} The vertex algebra $\mathcal{W}_{-1}(\mathfrak{sl}(4), f_{\rm sub})$ is the Kazama--Suzuki dual of the $N=2$ superconformal vertex algebra $L^{N=2}_{c=-15}$.
\end{Theorem}

\appendix\setcounter{section}{1}
% Original source lines 1108--1148
 \section[Operator product expansions for W\^{}k(sl\_4, f\_sub)]{Operator product expansions for $\boldsymbol{\mathcal{W}^{k}(\mathfrak{sl}_4, f_{\rm sub})}$}
 \label{appendix-b}

 The subregular $\mathcal{W}$-algebra $\mathcal{W}^{k}(\mathfrak{sl}_4, f_{\rm sub})$ is isomorphic to the Feigin--Semikhatov algebra $W^{(2)}_4$. It is generated by the even fields $J(z)$, $L(z)$, $G^{\pm}(z)$, $W(z)$ satisfying the following OPEs (cf.~\cite{FS}):\looseness=1
 \begin{gather*}
 J(x)J(y) \sim \frac{3k+8}{4(z-w)^{2}}, \qquad
 J(z)G^{\pm}(w) \sim \pm \frac{G^{\pm}(w)}{z-w} , \enskip \qquad G^{\pm}(z)G^{\pm}(w) \sim0, \\
	L(z)G^{\pm}(w) \sim \frac{2G^{\pm}(w)}{(z-w)^{2}} + \frac{\partial G^{\pm}(w)}{(z-w)}, \qquad
	L(z)J(w) \sim \frac{J(w)}{(z-w)^{2}} + \frac{\partial J(w)}{(z-w)}, \\
	L(z)L(w) \sim - \frac{c_k}{2(z-w)^{4}}+ \frac{2L(w)}{(z-w)^{2}}+ \frac{\partial L(w)}{z-w}, \\
	L(z) W(w) \sim \frac{3 W(w)}{(z-w)^{2}} + \frac{ \partial W(w)}{(z-w)}, \qquad J(z) W(w) \sim 0,
\\
	G^+(z)G^-(w) \sim \frac{(k+2)(2k+5)(3k+8)\mathbbm{1}}{(z-w)^{4}}+ \frac{4(k+2)(2k+5)J(w)}{ (z-w)^{3}} \\
	\phantom{	G^+(z)G^-(w) \sim }{} + (k+2) \frac{6 {:}JJ{:}(w) +2(2k+5)\partial J(w)-(k+4)L(w)}{(z-w)^{2}} \\
		\phantom{	G^+(z)G^-(w) \sim }{} + (k+2)\left( W(w) + \left(\frac{8(k+11)}{3(3k+8)^2}{:}J^3{:}(w)-\frac{4(k+4)}{3k+8}{:}L(w)J(w){:}\right.\right.\\
\left.\left.\phantom{	G^+(z)G^-(w) \sim }{}+ 6{:}\partial J(w)J(w){:}	 + \frac{k+4}{2}\partial L(w)\right.\right. \\
\left.\left.\phantom{	G^+(z)G^-(w) \sim }{}+ \frac{4(3k^2+17k+26)}{3(3k+8)}\partial ^2J(w)\right)\right)(z-w)^{-1},
\\
	 W(z)G^{\pm}(w) \sim \pm\frac{2(k+4)(3k+7)(5k+16)G^{\pm}(w)}{(3k+8)^2(z-w)^{3}}+ \left(\pm\frac{3(k+4)(5k+16)\partial G^{\pm}(w)}{2(3k+8)} \right. \\
\phantom{ W(z)G^{\pm}(w) \sim }{}	 \left. -\frac{6(k+4)(5k+16){:}J(w)G^{\pm}(w){:} }{(3k+8)^2}\right)(z-w)^{-2} \\
\phantom{ W(z)G^{\pm}(w) \sim }{}+\left(-\frac{8(k+4)(k+3){:}J(w)\partial G^{\pm}(w){:}}{(k+2)(3k+8)} \right. \\
\phantom{ W(z)G^{\pm}(w) \sim }{}	 \left. - \frac{4(k+4)(3k^2+15k+16){:}\partial J(w)G^{\pm}(w){:}}{(k+2)(3k+8)^2}\right.\\
 \phantom{ W(z)G^{\pm}(w) \sim }{}\left. \pm \frac{(k+4)(k+3)\partial ^2G^{\pm}(w)}{(k+2)}\mp \frac{2(k+4)^2{:}L(w)G^{\pm}(w){:}}{(k+2)(3k+8)}\right. \\
	 \phantom{ W(z)G^{\pm}(w) \sim }{} \left. \pm \frac{4(k+4)(5k+16){:}J(w)^2G^{\pm}(w){:}}{(k+2)(3k+8)^2} \right)(z-w)^{-1} , \\
	 W(z)W(w) \sim \frac{2(k+4)(2k+5)(3k+7)(5k+16)\mathbbm{1}}{(3k+8)(z-w)^{6}} -\frac{3(k + 4)^2( 5 k + 16)}{(3k+8)(z-w)^{4}}L^{\perp}(w) \\
\phantom{W(z)W(w) \sim}{}	 -\frac{3(k + 4)^2( 5 k + 16)}{2(3k+8)(z-w)^{3}}\partial L^{\perp}(w) \\
\phantom{W(z)W(w) \sim}{}+\left( -\frac{3(k + 4)^2( 5 k + 16)(12k^2 + 59k + 74)}{4(3k+8)(20k^2 + 93k + 102)}\partial^2L ^{\perp}(w)\right. \\
\phantom{W(z)W(w) \sim}{}	 + \left. \frac{8(k + 4)^3( 5 k + 16)}{(3k+8)(20k^2 + 93k + 102)}{:}L^{\perp}(w)L^{\perp}(w){:} + 4 (k+4) \Lambda (w)\right) (z-w)^{-2} \\
\phantom{W(z)W(w) \sim}{}	 + \biggl( -\frac{(k + 4)^2( 5 k + 16)(12k^2 + 59k + 74)}{6(3k+8)(20k^2 + 93k + 102)}\partial ^3L^{\perp}(w) \\
\phantom{W(z)W(w) \sim}{}	 +   \frac{8(k + 4)^3( 5 k + 16)}{(3k+8)(20k^2 + 93k + 102)}{:}\partial L^{\perp}(w)L^{\perp}(w){:} \\
\phantom{W(z)W(w) \sim}{}  + 2 (k+4) \partial \Lambda(w)\biggr)(z-w)^{-1},
\end{gather*}
 where
 \begin{align*}
 	(k+2)^2\Lambda = {}&{:} G^+ G^-{:} + (k+2) \left(-\frac{\partial W(z)}{2} -\frac{4 {:}W(z)J(z){:}}{3k+8}\right. \\
 &\left.+ \frac{(k+2)(k+4)\bigl(6k^2+33k+46\bigr)}{2(3k+8)\bigl(20k^2+93k+102\bigr)}\partial^2L^{\perp}(z)\right. \\
 	 &- \left. \frac{(k+4)^2(11k+26)}{2(3k+8)\bigl(20k^2+93k+102\bigr)}{:}L^{\perp}(z)L^{\perp}(z){:} + \frac{2(k+4)}{3k+8}\partial {:}L^{\perp}(z)J(z){:} \right.\\
 	 &+ \left. \frac{8(k+4)}{(3k+8)^2}{:}L^{\perp}(z)J(z)J(z){:} - \frac{2k+5}{3k+8}\left(\frac{8}{3}{:}\partial ^2J(z)J(z){:} + 2{:}\partial J(z)\partial J(z){:}\right. \right. \\
 	 &+ \left. \left. \frac{16}{3k+8}{:}\partial J(z)J(z)J(z){:} + \frac{32}{3(3k+8)^2}{:}J(z)^4{:} + \frac{3k+8}{6}\partial ^3J(z)\right) \right),
 \end{align*}
 $L^{\perp} = L-\frac{2}{3k+8}{:}JJ{:}$ and $c_k = -\frac{(3k+8)(8k+17)}{(k+4)}$.
\begin{thebibliography}{99}
\bibitem[1]{A-IMRN}
Adamovi\'c D., Representations of the~{$N=2$} superconformal vertex algebra,
 \href{https://doi.org/10.1155/S1073792899000033}{\textit{Int. Math. Res. Not.}} \textbf{1999} (1999), 61--79,
 \href{http://arxiv.org/abs/math.QA/9809141}{arXiv:math.QA/9809141}.

\bibitem[2]{A-2001}
Adamovi\'c D., Vertex algebra approach to fusion rules for~${N=2}$
 superconformal minimal models,
 \href{https://doi.org/10.1006/jabr.2000.8728}{\textit{J.~Algebra}}
 \textbf{239} (2001), 549--572.

\bibitem[3]{A-2005}
Adamovi\'c D., A construction of admissible {$A^{(1)}_1$}-modules of level
 {$-\frac 43$},
 \href{https://doi.org/10.1016/j.jpaa.2004.08.007}{\textit{J.~Pure Appl.
 Algebra}} \textbf{196} (2005), 119--134,
 \href{http://arxiv.org/abs/math.QA/0401023}{arXiv:math.QA/0401023}.

\bibitem[6]{AMW}
Adamovi\'c D., Milas A., Wang Q., On parafermion vertex algebras
 of~{$\mathfrak{sl}(2)$} and~{$\mathfrak{sl}(3)$} at level~{$-\frac32$},
 \href{https://doi.org/10.1142/S0219199720500868}{\textit{Commun. Contemp.
 Math.}} \textbf{24} (2022), 2050086, 23~paages,
 \href{http://arxiv.org/abs/2005.02631}{arXiv:2005.02631}.

\bibitem[10]{BEHHH}
Blumenhagen R., Eholzer W., Honecker A., H\"{u}bel R., Hornfeck K., Coset
 realization of unifying~{$\mathcal W$} algebras,
 \href{https://doi.org/10.1142/S0217751X95001157}{\textit{Internat.~J. Modern
 Phys.~A}} \textbf{10} (1995), 2367--2430,
 \href{http://arxiv.org/abs/hep-th/9406203}{arXiv:hep-th/9406203}.

\bibitem[13]{CL18}
Creutzig T., Linshaw A.R., Cosets of the {$\mathcal{W}^k(\mathfrak{sl}_4,
 f_{\rm subreg})$}-algebra, in Vertex {A}lgebras and {G}eometry,
 \textit{Contemp. Math.}, Vol. 711,
 \href{https://doi.org/10.1090/conm/711/14301}{American Mathematical Society},
 Providence, RI, 2018, 105--117,
 \href{http://arxiv.org/abs/1711.11109}{arXiv:1711.11109}.

\bibitem[15]{DLY}
Dong C., Lam C.H., Yamada H., {$\mathcal W$}-algebras related to parafermion
 algebras,
 \href{https://doi.org/10.1016/j.jalgebra.2009.03.034}{\textit{J.~Algebra}}
 \textbf{322} (2009), 2366--2403,
 \href{http://arxiv.org/abs/0809.3630}{arXiv:0809.3630}.

\bibitem[16]{Feh}
Fehily Z., Subregular {$\mathcal W$}-algebras of type~{$A$},
 \href{https://doi.org/10.1142/S0219199722500493}{\textit{Commun. Contemp.
 Math.}} \textbf{25} (2023), 2250049, 44~pages,
 \href{http://arxiv.org/abs/2111.05536}{arXiv:2111.05536}.

\bibitem[17]{FS}
Feigin B.L., Semikhatov A.M., {$\mathcal W^{(2)}_n$}~algebras,
 \href{https://doi.org/10.1016/j.nuclphysb.2004.06.056}{\textit{Nuclear
 Phys.~B}} \textbf{698} (2004), 409--449,
 \href{http://arxiv.org/abs/math.QA/0401164}{arXiv:math.QA/0401164}.

\bibitem[18]{FST}
Feigin B.L., Semikhatov A.M., Tipunin I.Y., Equivalence between chain
 categories of representations of affine~{$\mathfrak{sl}(2)$} and~{$N=2$}
 superconformal algebras,
 \href{https://doi.org/10.1063/1.532473}{\textit{J.~Math. Phys.}} \textbf{39}
 (1998), 3865--3905,
 \href{http://arxiv.org/abs/hep-th/9701043}{arXiv:hep-th/9701043}.

\bibitem[19]{Gen}
Genra N., Screening operators for {$\mathcal{W}$}-algebras,
 \href{https://doi.org/10.1007/s00029-017-0315-9}{\textit{Selecta Math.
 (N.S.)}} \textbf{23} (2017), 2157--2202,
 \href{http://arxiv.org/abs/1606.0096}{arXiv:1606.0096}.

\bibitem[21]{GK}
Gorelik M., Kac V., On simplicity of vacuum modules,
 \href{https://doi.org/10.1016/j.aim.2006.09.004}{\textit{Adv. Math.}}
 \textbf{211} (2007), 621--677,
 \href{http://arxiv.org/abs/math-ph/0606002}{arXiv:math-ph/0606002}.

\bibitem[23]{KW}
Kac V., Roan S.S., Wakimoto M., Quantum reduction for affine superalgebras,
 \href{https://doi.org/10.1007/s00220-003-0926-1}{\textit{Comm. Math. Phys.}}
 \textbf{241} (2003), 307--342,
 \href{http://arxiv.org/abs/math-ph/0302015}{arXiv:math-ph/0302015}.

\bibitem[25]{LL}
Lepowsky J., Li H., Introduction to vertex operator algebras and their
 representations, \textit{Progr. Math.}, Vol.~227,
 \href{https://doi.org/10.1007/978-0-8176-8186-9}{Birkh\"auser}, Boston, MA,
 2004.

\bibitem[26]{Li-ext}
Li H., Certain extensions of vertex operator algebras of affine type,
 \href{https://doi.org/10.1007/s002200100386}{\textit{Comm. Math. Phys.}}
 \textbf{217} (2001), 653--696,
 \href{http://arxiv.org/abs/math.QA/0003038}{arXiv:math.QA/0003038}.

\end{thebibliography}
\end{document}
